\begin{table}[H]
\centering
\caption{Forma de $y_{p}$ en $r(x)$ con forma trigónométrica}
\label{tab:eq_trigonométrica}
\begin{threeparttable}
{\setstretch{1.75}
\begin{tabular}{p{0.45\linewidth}p{0.45\linewidth}}\hline

    \textbf{Condición sobre $\pm \beta i$} & \textbf{Forma tentativa de $y_{p}$} \\ \hline

    $\pm \beta i$ no es raiz ($z = 0$) & $A \cos{(\beta x)} + B \sin{(\beta x)}$ \\
    $\pm \beta i$ si es raiz de multiplicidad $z$ & $x^{z} \cdot [A \cos{(\beta x)} + B \sin{(\beta x)}]$ \\
    
    \hline
\end{tabular}
}
\begin{tablenotes} % ex: \tntote{1} ... \item[1]
    \item[] 
\end{tablenotes}
\end{threeparttable}
\end{table}